\documentclass[11pt]{amsart}
\usepackage{geometry}                % See geometry.pdf to learn the layout options. There are lots.
\geometry{letterpaper}                   % ... or a4paper or a5paper or ... 
%\geometry{landscape}                % Activate for for rotated page geometry
%\usepackage[parfill]{parskip}    % Activate to begin paragraphs with an empty line rather than an indent
\usepackage{graphicx}
\usepackage{amssymb}
\usepackage{epstopdf}
\DeclareGraphicsRule{.tif}{png}{.png}{`convert #1 `dirname #1`/`basename #1 .tif`.png}

\title{Notes on solution to the buoyancy equation with exponential $\kappa$ profile}
\author{GPT 5.5 with  prompting from TWNH }
\date{6Jul26}                                           % Activate to display a given date or no date

\begin{document}
\maketitle

\section{Problem}

Consider the homogeneous ODE associated with the vertical Green's function problem
\begin{equation}
    \frac{Ek}{\Gamma}
    \frac{d}{dz}
    \left(
        \kappa(z)\frac{db}{dz}
    \right)
    - b = 0.
\end{equation}
Equivalently, define
\begin{equation}
    \mu^2 = \frac{\Gamma}{Ek},
\end{equation}
so that the ODE becomes
\begin{equation}
    \frac{d}{dz}
    \left(
        \kappa(z)b_z
    \right)
    - \mu^2 b = 0.
    \label{eq:homogeneous_ode}
\end{equation}

We consider the exponential diffusivity profile
\begin{equation}
    \kappa(z)
    =
    a + \exp\left[-\frac{z+H}{h}\right],
    \label{eq:kappa_profile}
\end{equation}
where, in the specific example of interest,
\begin{equation}
    a = 10^{-2},
    \qquad
    h = 0.1.
\end{equation}

Thus the ODE is
\begin{equation}
    \frac{d}{dz}
    \left[
        \left(
            a + \exp\left[-\frac{z+H}{h}\right]
        \right)b_z
    \right]
    - \mu^2 b = 0.
    \label{eq:ode_exp_kappa}
\end{equation}

\section{Change of Variables}

Let
\begin{equation}
    s = z+H,
\end{equation}
so that \(s=0\) at the bottom \(z=-H\) and \(s=H\) at the surface \(z=0\). Then
\begin{equation}
    \kappa(s) = a + e^{-s/h}.
\end{equation}

Now define a nondimensional exponential coordinate
\begin{equation}
    r
    =
    \frac{e^{-s/h}}{a}
    =
    \frac{1}{a}
    \exp\left[-\frac{z+H}{h}\right].
    \label{eq:r_definition}
\end{equation}
Then
\begin{equation}
    e^{-s/h} = ar,
\end{equation}
and therefore
\begin{equation}
    \kappa = a(1+r).
\end{equation}

Moreover,
\begin{equation}
    \frac{dr}{dz}
    =
    -\frac{1}{h}r,
\end{equation}
so that
\begin{equation}
    \frac{d}{dz}
    =
    -\frac{r}{h}\frac{d}{dr}.
    \label{eq:dz_to_dr}
\end{equation}

Let
\begin{equation}
    b(z)=B(r).
\end{equation}
Then
\begin{equation}
    b_z
    =
    \frac{dB}{dr}\frac{dr}{dz}
    =
    -\frac{r}{h}B_r.
\end{equation}
Therefore,
\begin{equation}
    \kappa b_z
    =
    a(1+r)
    \left(
        -\frac{r}{h}B_r
    \right)
    =
    -\frac{a}{h}r(1+r)B_r.
\end{equation}

Differentiating once more,
\begin{align}
    \frac{d}{dz}
    \left(
        \kappa b_z
    \right)
    &=
    -\frac{r}{h}
    \frac{d}{dr}
    \left[
        -\frac{a}{h}r(1+r)B_r
    \right] \\
    &=
    \frac{ar}{h^2}
    \frac{d}{dr}
    \left[
        r(1+r)B_r
    \right].
\end{align}

Since
\begin{equation}
    \frac{d}{dr}
    \left[
        r(1+r)B_r
    \right]
    =
    r(1+r)B_{rr}
    +
    (1+2r)B_r,
\end{equation}
equation \eqref{eq:ode_exp_kappa} becomes
\begin{equation}
    \frac{ar}{h^2}
    \left[
        r(1+r)B_{rr}
        +
        (1+2r)B_r
    \right]
    -
    \mu^2 B
    =
    0.
\end{equation}

Dividing by \(a/h^2\) gives
\begin{equation}
    r
    \left[
        r(1+r)B_{rr}
        +
        (1+2r)B_r
    \right]
    -
    \frac{\mu^2 h^2}{a}B
    =
    0.
\end{equation}
Hence
\begin{equation}
    r^2(1+r)B_{rr}
    +
    r(1+2r)B_r
    -
    \beta^2 B
    =
    0,
    \label{eq:B_equation}
\end{equation}
where
\begin{equation}
    \beta^2
    =
    \frac{\mu^2h^2}{a}
    =
    \frac{\Gamma h^2}{a\,Ek}.
    \label{eq:beta_squared}
\end{equation}
Thus
\begin{equation}
    \beta
    =
    h\sqrt{\frac{\Gamma}{a\,Ek}}.
    \label{eq:beta}
\end{equation}

\section{Reduction to Hypergeometric Form}

Equation \eqref{eq:B_equation} has a regular singular point at \(r=0\).
To determine the leading-order behavior, try
\begin{equation}
    B(r) = r^\lambda.
\end{equation}
Near \(r=0\), equation \eqref{eq:B_equation} reduces to
\begin{equation}
    r^2B_{rr} + rB_r - \beta^2B = 0.
\end{equation}
Substituting \(B=r^\lambda\) gives
\begin{equation}
    \lambda(\lambda-1)+\lambda-\beta^2=0,
\end{equation}
or
\begin{equation}
    \lambda^2-\beta^2=0.
\end{equation}
Therefore,
\begin{equation}
    \lambda=\pm\beta.
\end{equation}

Choose the positive exponent and write
\begin{equation}
    B(r)=r^\beta Y(r).
    \label{eq:B_substitution}
\end{equation}
Then
\begin{equation}
    B_r
    =
    \beta r^{\beta-1}Y
    +
    r^\beta Y_r,
\end{equation}
and
\begin{equation}
    B_{rr}
    =
    \beta(\beta-1)r^{\beta-2}Y
    +
    2\beta r^{\beta-1}Y_r
    +
    r^\beta Y_{rr}.
\end{equation}

Substitution into equation \eqref{eq:B_equation}, followed by simplification, yields
\begin{equation}
    r(1+r)Y_{rr}
    +
    \left[
        2\beta+1+(2\beta+2)r
    \right]Y_r
    +
    \beta(\beta+1)Y
    =
    0.
    \label{eq:Y_r_equation}
\end{equation}

This equation is transformed to the standard Gauss hypergeometric equation by setting
\begin{equation}
    x=-r.
    \label{eq:x_definition}
\end{equation}
Since
\begin{equation}
    \frac{dY}{dr}
    =
    -\frac{dY}{dx},
    \qquad
    \frac{d^2Y}{dr^2}
    =
    \frac{d^2Y}{dx^2},
\end{equation}
equation \eqref{eq:Y_r_equation} becomes
\begin{equation}
    x(1-x)Y_{xx}
    +
    \left[
        2\beta+1-(2\beta+2)x
    \right]Y_x
    -
    \beta(\beta+1)Y
    =
    0.
    \label{eq:hypergeometric_form}
\end{equation}

The standard Gauss hypergeometric equation is
\begin{equation}
    x(1-x)Y_{xx}
    +
    \left[
        c-(A+B+1)x
    \right]Y_x
    -
    ABY
    =
    0.
    \label{eq:standard_hypergeometric}
\end{equation}
Comparison of \eqref{eq:hypergeometric_form} and
\eqref{eq:standard_hypergeometric} gives
\begin{equation}
    c = 2\beta+1,
\end{equation}
\begin{equation}
    A+B+1 = 2\beta+2,
\end{equation}
and
\begin{equation}
    AB = \beta(\beta+1).
\end{equation}
One convenient choice is
\begin{equation}
    A=\beta,
    \qquad
    B=\beta+1,
    \qquad
    c=2\beta+1.
    \label{eq:hypergeometric_parameters}
\end{equation}

\section{Analytic Homogeneous Solutions}

A standard solution of the Gauss hypergeometric equation is
\begin{equation}
    Y_1(x)
    =
    {}_2F_1(A,B;c;x).
\end{equation}
Using \eqref{eq:hypergeometric_parameters},
\begin{equation}
    Y_1(x)
    =
    {}_2F_1
    \left(
        \beta,\beta+1;2\beta+1;x
    \right).
\end{equation}

Since \(x=-r\) and \(B(r)=r^\beta Y(r)\), one homogeneous solution is
\begin{equation}
    B_1(r)
    =
    r^\beta
    {}_2F_1
    \left(
        \beta,\beta+1;2\beta+1;-r
    \right).
\end{equation}

Returning to the original coordinate \(z\), with
\begin{equation}
    r(z)
    =
    \frac{1}{a}
    \exp\left[-\frac{z+H}{h}\right],
\end{equation}
the first solution is
\begin{equation}
    b_1(z)
    =
    r(z)^\beta
    {}_2F_1
    \left(
        \beta,\beta+1;2\beta+1;-r(z)
    \right).
    \label{eq:b1_solution}
\end{equation}

The second linearly independent solution corresponds to the other indicial exponent \(-\beta\). It is
\begin{equation}
    b_2(z)
    =
    r(z)^{-\beta}
    {}_2F_1
    \left(
        -\beta,1-\beta;1-2\beta;-r(z)
    \right).
    \label{eq:b2_solution}
\end{equation}

Therefore, the general homogeneous solution is
\begin{align}
    b(z)
    &=
    C_1
    r(z)^\beta
    {}_2F_1
    \left(
        \beta,\beta+1;2\beta+1;-r(z)
    \right)
    \nonumber\\
    &\quad
    +
    C_2
    r(z)^{-\beta}
    {}_2F_1
    \left(
        -\beta,1-\beta;1-2\beta;-r(z)
    \right).
    \label{eq:general_solution}
\end{align}

\section{Specialization to \(a=10^{-2}\), \(h=0.1\)}

For the specific diffusivity profile
\begin{equation}
    \kappa(z)
    =
    10^{-2}
    +
    \exp\left[-\frac{z+H}{0.1}\right],
\end{equation}
we have
\begin{equation}
    a=10^{-2},
    \qquad
    h=0.1.
\end{equation}
Then
\begin{equation}
    r(z)
    =
    \frac{1}{10^{-2}}
    \exp\left[-\frac{z+H}{0.1}\right]
    =
    100\exp[-10(z+H)].
\end{equation}
Furthermore,
\begin{equation}
    \beta
    =
    h\sqrt{\frac{\Gamma}{a\,Ek}}
    =
    0.1
    \sqrt{
        \frac{\Gamma}{10^{-2}Ek}
    }
    =
    \sqrt{\frac{\Gamma}{Ek}}.
\end{equation}
Thus, for this profile,
\begin{equation}
    \beta=\sqrt{\frac{\Gamma}{Ek}}.
\end{equation}

The homogeneous solution becomes
\begin{align}
    b(z)
    &=
    C_1
    r(z)^\beta
    {}_2F_1
    \left(
        \beta,\beta+1;2\beta+1;-r(z)
    \right)
    \nonumber\\
    &\quad
    +
    C_2
    r(z)^{-\beta}
    {}_2F_1
    \left(
        -\beta,1-\beta;1-2\beta;-r(z)
    \right),
\end{align}
with
\begin{equation}
    r(z)=100\exp[-10(z+H)],
    \qquad
    \beta=\sqrt{\frac{\Gamma}{Ek}}.
\end{equation}

\section{Green's Function Construction}

The Green's function \(G_b(z,\xi)\) satisfies
\begin{equation}
    \frac{Ek}{\Gamma}
    \frac{d}{dz}
    \left(
        \kappa(z)\frac{\partial G_b}{\partial z}
    \right)
    -
    G_b
    =
    -\delta(z-\xi),
    \label{eq:green_problem}
\end{equation}
with boundary conditions
\begin{equation}
    \frac{\partial G_b}{\partial z}(-H,\xi)=0,
    \qquad
    G_b(0,\xi)=0.
\end{equation}

Let \(u(z)\) be the homogeneous solution satisfying the bottom Neumann condition
\begin{equation}
    u'(-H)=0,
\end{equation}
and let \(v(z)\) be the homogeneous solution satisfying the surface Dirichlet condition
\begin{equation}
    v(0)=0.
\end{equation}

Both \(u\) and \(v\) are linear combinations of the two basis solutions
\eqref{eq:b1_solution} and \eqref{eq:b2_solution}. The generalized Wronskian
\begin{equation}
    W
    =
    \kappa(z)
    \left[
        u(z)v'(z)
        -
        u'(z)v(z)
    \right]
    \label{eq:wronskian}
\end{equation}
is independent of \(z\).

The Green's function is then
\begin{equation}
    G_b(z,\xi)
    =
    -\frac{\Gamma}{Ek\,W}
    \begin{cases}
        u(z)v(\xi), & z\le \xi, \\[2mm]
        u(\xi)v(z), & z\ge \xi.
    \end{cases}
    \label{eq:green_function}
\end{equation}

The jump condition follows by integrating \eqref{eq:green_problem}
across \(z=\xi\):
\begin{equation}
    \frac{Ek}{\Gamma}
    \left[
        \kappa G_{b,z}
    \right]_{\xi^-}^{\xi^+}
    =
    -1,
\end{equation}
or
\begin{equation}
    \left[
        \kappa G_{b,z}
    \right]_{\xi^-}^{\xi^+}
    =
    -\frac{\Gamma}{Ek}.
    \label{eq:jump_condition}
\end{equation}

\section{Summary}

Starting from
\begin{equation}
    \frac{d}{dz}
    \left[
        \left(
            a+e^{-(z+H)/h}
        \right)b_z
    \right]
    -
    \frac{\Gamma}{Ek}b
    =
    0,
\end{equation}
define
\begin{equation}
    r
    =
    \frac{1}{a}e^{-(z+H)/h},
    \qquad
    x=-r,
    \qquad
    \beta
    =
    h\sqrt{\frac{\Gamma}{a\,Ek}}.
\end{equation}
Then, with
\begin{equation}
    b(z)=r^\beta Y(x),
\end{equation}
the ODE becomes the standard Gauss hypergeometric equation
\begin{equation}
    x(1-x)Y_{xx}
    +
    \left[
        2\beta+1-(2\beta+2)x
    \right]Y_x
    -
    \beta(\beta+1)Y
    =
    0,
\end{equation}
corresponding to parameters
\begin{equation}
    A=\beta,
    \qquad
    B=\beta+1,
    \qquad
    c=2\beta+1.
\end{equation}

Therefore the homogeneous solution is
\begin{align}
    b(z)
    &=
    C_1
    r(z)^\beta
    {}_2F_1
    \left(
        \beta,\beta+1;2\beta+1;-r(z)
    \right)
    \nonumber\\
    &\quad
    +
    C_2
    r(z)^{-\beta}
    {}_2F_1
    \left(
        -\beta,1-\beta;1-2\beta;-r(z)
    \right).
\end{align}

\begin{thebibliography}{9}

\bibitem{abramowitz_stegun}
M.~Abramowitz and I.~A. Stegun,
\textit{Handbook of Mathematical Functions with Formulas, Graphs, and Mathematical Tables},
Dover Publications, 1965.

\bibitem{dlmf}
NIST Digital Library of Mathematical Functions,
\textit{Chapter 15: Hypergeometric Function},
\url{https://dlmf.nist.gov/15}.

\bibitem{olver}
F.~W.~J. Olver, D.~W. Lozier, R.~F. Boisvert, and C.~W. Clark, editors,
\textit{NIST Handbook of Mathematical Functions},
Cambridge University Press, 2010.

\bibitem{arfken}
G.~B. Arfken, H.~J. Weber, and F.~E. Harris,
\textit{Mathematical Methods for Physicists},
7th ed.,
Academic Press, 2013.

\bibitem{courant_hilbert}
R.~Courant and D.~Hilbert,
\textit{Methods of Mathematical Physics, Volume I},
Wiley-Interscience, 1989.

\bibitem{zauderer}
E.~Zauderer,
\textit{Partial Differential Equations of Applied Mathematics},
3rd ed.,
Wiley-Interscience, 2006.

\end{thebibliography}

\end{document}  